\chapter{Notation}
\section{Special operators and symbols}
\subsection*{Matrix diagonal extraction}
Operator $\diag$ extracts the elements on the
main diagonal of a matrix, i.e.\
$\diag\T{A} = a_{ii}$.

\subsection*{Matrix direct sum}
A direct sum of two matrices is defined as
\begin{equation}
  \T{A} \oplus \T{B} =
  \begin{bmatrix}
    \T{A} & \T{0} \\
    \T{0} & \T{B}
  \end{bmatrix}.
\end{equation}
Furthermore, a sum of $n$ elements is defined as
\begin{equation}
  \bigoplus_{i = 1}^n \T{A}_i =
  \begin{bmatrix}
    \T{A}_1 & \T{0}   & \cdots & \T{0}   \\
    \T{0}   & \T{A}_2 & \cdots & \T{0}   \\
    \vdots  & \vdots  & \ddots & \vdots  \\
    \T{0}   & \T{0}   & \cdots & \T{A}_n
  \end{bmatrix}.
\end{equation}

\subsection*{Inner product}
A bilinear form $\iprd[\T{A}]{\,\cdot\,}{\,\cdot\,}  : \Complex^n \times \Complex^n \to \Complex$
named $\T{A}$-inner product is defined as $\iprd[\T{A}]{\T{u}}{\T{v}} = \T{u}^* \T{A} \T{v}$.
In the special case of $\T{A} = \T{I}$, which is denoted by simply
$\iprd{\,\cdot\,}{\,\cdot\,}$, the
inner product is equivalent to the dot product of two vectors:
$\iprd[\T{I}]{\T{u}}{\T{v}} = \iprd{\T{u}}{\T{v}} = \T{u}^* \T{v}$.

\section{\nameCaSymbols}

\subsubsection*{Typographical symbols}
\S\ chapter, section, subsection

$\blacksquare$ end of a proof

$\vartriangle$ end of a remark or example

$\blacktriangle$ end of a definition

\subsubsection*{Mathematical symbols}

$j$ imaginary unit

$\T{I}$ identity matrix

$\T{0}$ zero matrix

$\T{P}$ permutation matrix

$\T{Q}$, $\T{U}$ unitary (orthogonal) matrix

$\T{T}$ transformation matrix

$\T{\mathit{R}}(\phi)$ rotation matrix

$\cond{(\T{A})}$ condition number of $\T{A}$

$\lambda(\T{A})$ eigenvalue of $\T{A}$

$\sigma(\T{A})$ singular value of $\T{A}$

$\mathscr{H}_p$ Hardy $p$-norm

$\mathscr{L}_p$ Lebesgue $p$-norm

\section{\nameCaAcronyms}
\label{sc:abbreviationsAndAcronyms}
%\vspace{-1 cm}
%\printacronyms[name = {}]
\printacronyms[name={}, heading=none]
